\section{Search strategy and screening pipeline}
\label{app:search}

\subsection{Database, scope and time window}
\label{app:search:scope}

The literature search was conducted in Scopus advanced search over
the time window 2016--2026 inclusive. Document types are restricted
to article, review and conference paper. The language is English.
The subject areas are energy, engineering, computer science,
environmental science and mathematics. The year filter is set in the
Scopus user interface; the query string itself contains neither
\texttt{TITLE-ABS-KEY} nor \texttt{PUBYEAR} so that the same query
body can be reused with adjusted side-bar filters.

\subsection{Active search strings}
\label{app:search:strings}

Two complementary search strings define the union from which the
evidence base is built. The first targets the explicit surrogate-
modelling vocabulary in conjunction with energy-system and
optimization context. The second imports the user's earlier
multi-criteria multi-objective Scopus export, which serves as the
application backbone for multi-energy systems, microgrids, energy
communities, hybrid renewable energy systems, district heating and
cooling, combined cooling-heating-power and sector-coupled systems.
Two narrower helper queries (surrogate~$\times$~stochastic/robust
and surrogate~$\times$~multi-objective) ensure that nuanced
sub-fields are not lost in the noise of the main search.

\subsection{Offline screening pipeline}
\label{app:search:offline}

The Scopus export of the main search returned 3129 entries; the
multi-objective multicriteria export contributed 2369 additional
entries. A manual review of every record is not feasible, but a
fully automated keyword filter would either over-trim relevant
implicit surrogate work or under-trim noise. We therefore use a
tiered offline screening implemented in \texttt{filter\_bib.py} and
documented in \texttt{references/screening\_log.md}.

\paragraph{Tier A -- explicit surrogate work.} Entries match a closed
vocabulary of explicit surrogate-method terms (e.g. \emph{surrogate},
\emph{metamodel}, \emph{response surface}, \emph{kriging}, \emph{Gaussian
process}, \emph{polynomial chaos}, \emph{learning-to-optimize}) in
conjunction with a closed vocabulary of energy-system terms (e.g.
\emph{energy system}, \emph{power system}, \emph{district heating},
\emph{unit commitment}, \emph{optimal power flow}, \emph{microgrid}).
Tier-A entries are accepted automatically.

\paragraph{Tier B -- implicit surrogate work.} Entries that do not match
Tier A are accepted as candidates if they jointly satisfy four
conditions: (i) they mention a machine-learning or regression method,
(ii) they mention an optimization context, (iii) they mention an
energy-system context, and (iv) at least one ``proxy hint''
(\emph{proxy}, \emph{approximate}, \emph{accelerate}, \emph{replace},
\emph{computationally expensive}, etc.) signals that the model is used
as a substitute for an expensive component. Tier-B entries are
flagged for manual review.

\paragraph{Hard exclusions.} A noise vocabulary (\emph{fault diagnosis},
\emph{site selection}, \emph{image recognition}, \emph{drug discovery},
etc.) demotes Tier-B entries that show two or more strong noise
signals; Tier-A entries are flagged but kept, on the assumption that
explicit surrogate vocabulary outweighs an isolated noise term.

\paragraph{Merging.} The two raw exports are merged with
\texttt{build\_review\_bibliography.py}. Deduplication is DOI-first
and citation-key-second. Duplicate citation keys in the merged
bibliography are deterministically suffixed (e.g. \texttt{Wang2025}
$\rightarrow$ \texttt{Wang2025\_\_2}) so that all keys remain unique.
Per-entry provenance is recorded in
\texttt{review\_mes\_moo\_surrogates\_manifest.csv}.

\paragraph{Curated paper library.} A second screening stage
(\texttt{paper\_library/select\_paper\_library.py}) selects a curated
subset of approximately 260 entries that align with the manuscript
narrative. The selection enforces two conditions: every citation key
that appears in the manuscript text is included as a mandatory entry,
and the remaining slots are filled by a bucket-based top-up that
mirrors the manuscript's section structure (Sections~\ref{sec:taxonomy}
through~\ref{sec:applications}). The selection writes
\texttt{review\_paper\_library.bib}, a per-entry manifest, a buckets
table and a citation plan. Scopus-specific telemetry (URLs to
\texttt{scopus.com}, ``Cited by'' notes, source/publication-stage/
type fields, author keywords) is stripped on emission so the curated
bibliography is publication-grade. The screening pipeline is
deterministic and re-runnable from the raw exports.

\subsection{Run-level statistics}
\label{app:search:stats}

The combined manuscript bibliography
(\texttt{review\_mes\_moo\_surrogates.bib}) contains 2906 deduplicated
entries from 2944 inputs. The curated paper library
(\texttt{review\_paper\_library.bib}) contains 260 entries, of which
132 are mandatory citations from the manuscript text and 128 are
top-up entries selected by citation count within their assigned
bucket. The bucket-level distribution and per-entry mapping are
recorded in
\texttt{paper\_library/review\_paper\_library\_buckets.csv} and in
\texttt{paper\_library/review\_paper\_library\_citation\_plan.md};
they are reproduced in compact form in
Section~\ref{app:evidence-map}.

\section{Evidence map (extended manifest)}
\label{app:evidence-map}

The compact per-study evidence map in
Table~\ref{tab:T6-evidence-map} (Section~\ref{sec:applications})
lists the MES-focused subset used in the synthesis. The full curated
library cross-tabulates all 285 screened entries against the same
surrogate-family, task-class and integration-role axes in
\texttt{paper\_library/review\_paper\_library\_manifest.csv} and
\texttt{paper\_library/review\_paper\_library\_buckets.csv}. The
map is regenerated deterministically by
\texttt{tables/build\_table\_T6\_evidence\_map.py}.
